\documentclass[11pt]{article}
\usepackage[utf8]{inputenc}
\usepackage[T1]{fontenc}
\usepackage[italian]{babel}
\usepackage{geometry}
\geometry{a4paper, margin=1in}
\usepackage{xcolor}
\usepackage{listings}
\usepackage{hyperref}

\lstset{
    language=C,
    basicstyle=\ttfamily\small,
    keywordstyle=\color{blue},
    stringstyle=\color{red},
    commentstyle=\color{green!60!black},
    numbers=none,
    frame=single,
    breaklines=true
}

\usepackage{tcolorbox}
\tcbuselibrary{breakable}
\begin{document}

\begin{tcolorbox}[colback=blue!5!white,colframe=blue!75!black,title=\textbf{DOMANDA}, breakable]
7) Modificare il server in maniera che soddisfi 5 richieste prima di terminare
\begin{itemize}
    \item E se volessi che non terminasse mai?
\end{itemize}
\end{tcolorbox}

\begin{tcolorbox}[colback=green!5!white,colframe=green!75!black,title=\textbf{RISPOSTA}, breakable]
Ecco le risposte alle tue domande:

\begin{itemize}
    \item \textbf{Modificare il server per soddisfare 5 richieste prima di terminare:} Nel file \texttt{serverUDP.c}, ho sostituito il ciclo di ascolto infinito \texttt{while(true)} con un ciclo definito che itera esattamente 5 volte. Il codice del blocco di ascolto è diventato il seguente:
\end{itemize}

\begin{lstlisting}
// Al posto di while(true)
for(int i = 0; i < 5; i++) {
    printf("[SERVER] Sono in attesa di richieste da qualche client\n");
    UDPReceive(socket, request, MTU, hostAddress, &port);
    // ... stampe ...
    UDPSend(socket, response, strlen(response), hostAddress, port);
}
// Dopo aver servito i 5 client, il programma esce dal loop e il main() termina.
\end{lstlisting}

\begin{itemize}
    \item \textbf{E se volessi che non terminasse mai?} Per avere un server che resta in ascolto perenne senza terminare mai (comportamento tipico dei server veri e propri, detti anche daemoni), si utilizza un ciclo infinito. Nel linguaggio C, questo si ottiene utilizzando costrutti condizionali perennemente verificati, come:
    \begin{itemize}
        \item \texttt{while(true)} (utilizzato nella versione originale, richiede la libreria \texttt{<stdbool.h>})
        \item \texttt{while(1)} (costrutto tipico in C per indicare che la condizione non sarà mai 0)
        \item \texttt{for(;;)} (ciclo for privo di condizioni di terminazione)
    \end{itemize}
    In questo modo il programma attenderà all'infinito un nuovo pacchetto sulla porta senza fermarsi mai, a meno che non venga interrotto forzatamente (es. chiudendo il terminale o premendo \texttt{Ctrl+C}).
\end{itemize}
\end{tcolorbox}

\end{document}
